\bookchapter{Annexe A : Flux « niveau nounou » pour le déballage d'un nouveau smartphone}

\booksection{Les « trois étapes » du déballage et de la vérification :}

\begin{itemize}
  \item \textbf{Vérifier l’aspect extérieur}:Vérifiez si le châssis et l’écran présentent des rayures ou des chocs.
  \item \textbf{Vérifier les informations}:Vérifiez que l’IMEI/numéro de série sur la boîte correspond à celui sur la facture.
  \item \textbf{Tester les fonctions}:Passez un appel pour tester l’écouteur et le micro, prenez une photo pour vérifier la caméra.
\end{itemize}

\booksection{Activation et configuration initiale — « rappels importants » :}

\begin{itemize}
  \item \textbf{Réseau}:Connectez d'abord un Wi-Fi stable.
  \item \textbf{Compte}:Connectez votre Apple ID, compte Huawei, compte Xiaomi, etc.
  \item \textbf{Réglages de base}:Activez « Trouver mon smartphone » et définissez un mot de passe d'écran de verrouillage fiable.
  \item \textbf{Configuration eSIM}:Si votre smartphone prend en charge l'eSIM (comme l'iPhone et certains modèles phares Android), vous pouvez ajouter un forfait eSIM dans les paramètres, sans insérer de carte SIM physique. À l'étranger, l'eSIM est particulièrement pratique : achetez directement un forfait eSIM de l'opérateur local.
  \item \textbf{Migration des données}:Lors de la migration des données depuis l'ancien smartphone, assurez-vous que les deux appareils sont suffisamment chargés. iPhone : utilisez « Démarrage rapide » ; Android : utilisez l'appli officielle de transfert.
\end{itemize}

\booksection{Installation d'applications — « la sécurité d'abord » :}

\begin{itemize}
  \item \textbf{Privilégier le magasin officiel}:Apple : App Store ; Android : magasin d'applications intégré au smartphone.
  \item \textbf{Ne donnez pas les permissions au hasard}:Pour les demandes de permissions lors de l'installation, demandez-vous « cette appli en a-t-elle vraiment besoin ? »
\end{itemize}
